\section{光电效应}
在光的照射下，金属中的电子会吸收光能二逸出金属表面，这种现象称为\uwave{光电效应}，在光电效应中逸出金属表面的电子称为\uwave{光电子}，光电效应的实用器件是\uwave{光电管}，光电管是一个抽成真空的玻璃泡，如\xref{fig:光电管}所示，内表面的一部分涂有感光金属作为阴极，而阴极金属将在光照下产生光电子，光电子在外加电场的作用下向右作定向运动，被另一侧由金属丝网构成的阳极捕获，在电路中形成稳定的电流。我们将光电子在电场的作用下运动产生的电流，称为\uwave{光电流}。
\begin{Figure}[光电管]
    \includegraphics[scale=1]{build/Chapter05A_01.fig.pdf}
\end{Figure}

为了研究光电效应的规律，我们可以将光电管接入一个电压可变的电压源上，并通过电流表测定其产生的光电流。而实验表明，光电流的大小和电源电压之间存在如\xref{fig:光电流与电压的关系}所示的关系。

\subsection{光电效应的实验规律}

当施加正向电压时，光电流会随着电压的增大而增大，最终到达饱和，这很容易理解，因为光电效应产生的光电子的初始动能有大有小，有些初动能较大的电子不需要外加电压就可以到达阳极，有些初动能较小的电子则无法做到这一点，但是如果外加一些电压，那么这些电子也可以到达阳极变为光电流的一部分了。而之所以会到达饱和，是因为光强是一定的，当电压将足够大时，\empx{光电子逸出多少，光电流就有多少}，无论光电子的动能高低都能在强电场的作用下到达阳极参与进光电流，由于单位时间内产生的光电子的数目是由光强决定的，因此达到饱和后进一步增大电压并不会增大光电流，也因此，\empx{光强越大，光电流的饱和值就相应越高}。

当施加反向电压时，光电流不会立即为零，这是因为光电子是带有一定初动能，如果携带的动能足够，即便是面对起到阻碍作用的反向电场，其也能成功抵达阳极，而只有反向电压达到一定值时，光电流才能刚好为零，光电流刚好为零时光电管两端的电势差$U_0$称为\uwave{遏制电压}。

\begin{Figure}[光电流与电压的关系]
    \includegraphics[scale=0.75]{build/Chapter05A_02.fig.pdf}
\end{Figure}

很明显，光电子的最大初动能就是遏制电压消耗的电势能
\begin{Equation}
    \frac{1}{2}m_ev_{\text{m}}^2=e U_0
\end{Equation}
因此，如果我们希望研究光电子的最大初动能，就可以通过遏制电压的大小来研究。

\subsection{光电效应的诘难}
光电效应和光的波动说的矛盾究竟在哪里？问题就出在遏制电压，实验表明
\begin{Equation}
    U_0=k(\nu-\nu_0)
\end{Equation}
即，遏制电压$U_0$和$\nu$构成线性关系，其中$k$是普适的常数，而$\nu_0>0$则是某个与金属材料属性有关的常数。然而这个结果是非常古怪的，它给出的结果是，光电子的初始动能仅与光的频率有关，光强只会影响光电子的数目或光电流的强度，与光电子的初始动能无关。如果我们把矛盾显得再尖锐一些，这就是说，如果光的频率低于某个值，确切的说就是低于$\nu_0$，此时无论光强多大，都不会有任何光电子逸出，因为此时每个光电子的初始动能都是负值，它们无法脱离金属。我们将$\nu_0$称为金属的\uwave{截至频率}或\uwave{红限频率}。此处的红限与红光无关，只是因为在可见光波段，红光代表长波低频，紫光代表短波高频，我们只是借用了红对低频的指称。

光电效应中，截至频率的存在对于光的波动学说是不可接受的，光如果是一种波动，那么无论频率多少，只要光强足够大或光的照射持续足够长的时间，光电子总是能取得足够的初动能的。但实验事实是，如果光的频率低于截至频率，即便光强很大，也始终无法产生光电子。

光电效应对光的波动学说的另一项挑战是，如果光是一种波动，电子从光中吸收能量需要一定的时间，尤其是当光强较弱是，但实验结果却表明，光强的强弱与光电子产生的弛豫时间无关，而且即便光强很弱，光电子产生的弛豫时间仍然非常短，远低于光的波动说的预期值。

\subsection{爱因斯坦光电效应方程}
面对光电效应对光的波动学说提出的诘难，爱因斯坦在1905年结合光电效应的实验事实和普朗克于1900年在黑体辐射中提出的量子化假设，提出了下述的光子假说，\empx{光是一系列离散的以光速运动的光量子组成的粒子流}，换言之，光的能量是非连续的，是一份一份的，\uwave{光量子}就是光的能量的最小单位，即一份光的能量就是一个光量子。有时也将\uwave{光子}作为光量子的简称。

那么，光量子或者说光子具有多少能量呢？
\begin{BoxFormula}[光子的能量]
    光子的能量仅于光的频率有关
    \begin{Equation}&[]
        E=h\nu
    \end{Equation}
    其中$h$称为\uwave{普朗克常数}，其取值为$h=6.626\times 10^{-34}\si{J\cdot s}$。
\end{BoxFormula}
\empx{光子的能量与光的频率成正比}，而在一定频率下，光的强度，或者说，光的总能量其实就是在数有多少个光子，\empx{光的强度越大则光子的数目越多}，而单个光子的能量则是完全确定的。

那么，光量子假说下，光电效应将具有何种图景？由于电子在原子核外的广袤空间中是非常小的，实际上，电子与光子发生碰撞的概率是极小的，因此，我们可以作出这样的合理假设，即在光电效应中，\empx{电子实际上只有机会与一个光子发生碰撞}。电子和光子碰撞后，电子会在一瞬间吸收光子而获得一份能量$h\nu$，这部分能量首先将会消耗在电子逸出金属表面的过程中克服原子核的库仑力所必须做的逸出功$W$，当然，逸出功$W$有一个最低值，代表电子以最短的路径脱离原子，假如电子最初的运动方向不是很恰当，斜着甚至反向，那么所需的逸出功就会相应增加甚至超过其获得的能量，这样，电子的初动能就会相应减少甚至无法逃逸。而如果在最好的情况下，电子以最低的逸出功脱离了原子，那么此时，电子的初始动能就是最大初动能。
\begin{BoxEquation}[爱因斯坦光电效应方程]
    光电子吸收的能量，一部分转化为逸出功，一部分转化为初动能
    \begin{Equation}
        h\nu=\frac{1}{2}m_ev_{\text{m}}^2+W
    \end{Equation}
    其中$\nu$为入射光的频率，而$m_e$为电子质量。
\end{BoxEquation}\goodbreak

由此，光电效应的结果就很容易解释了，电子吸收光子是瞬时的，因此弛豫时间很短，光子的能量仅与频率无关，电子的能量仅能来自单个光子，电子逸出的关键在于其获得的能量是否达到最小逸出功，因此只有入射光的频率超过截至频率才有可能使得电子逸出，形成光电流。

现在，让我们来确定$U_0=k(\nu-\nu_0)$中的$k$和$\nu_0$吧。
\begin{BoxFormula}[遏制电压的关系式]
    光电效应的遏制电压与入射光频率具有线性关系
    \begin{Equation}
        U=k(\nu-\nu_0)
    \end{Equation}
    这种关系可以具体表示为
    \begin{Equation}
        U=\frac{h}{e}\qty(\nu-\frac{W}{h})
    \end{Equation}
\end{BoxFormula}
\begin{Proof}
    这很容易证明，遏制电压消耗的电势能就是光电子的最大初动能
    \begin{Equation}&[1]
        \frac{1}{2}m_ev_{\text{m}}^2=eU_0
    \end{Equation}
    即有
    \begin{Equation}&[2]
        U_0=\frac{1}{e}\qty(\frac{1}{2}m_{e}v_{\text{m}}^2)
    \end{Equation}
    运用\fancyref{eqt:爱因斯坦光电效应方程}
    \begin{Equation}
        U_0=\frac{1}{e}\qty(h\nu-W)
    \end{Equation}
    不妨将括号中的$h$提出
    \begin{Equation}*
        U_0=\frac{h}{e}\qty(\nu-\frac{W}{h})\qedhere
    \end{Equation}
\end{Proof}
爱因斯坦因为其对光电效应的理论解释，获得了1921年诺贝尔物理学奖，而两年之后，密里根由于对基本电荷（油滴实验）和普朗克常数的实验测定，获得了1923年诺贝尔物理学奖。

\subsection{光的波粒二象性}
现在我们看到，光的干涉和衍射证实了光的波动特性，光电效应证实了光的粒子特性，那么现在的问题是，光到底是波还是粒子？这两者似乎是不可调和的，波就是波，粒子就是粒子，而同一个事物怎么可能具有两种截然不同的存在。但这只是经典理论的看法在，我们不能将这种粗糙的直觉代入到微观世界的研究中，我们的实验其实只是表明，光具有波动性和粒子性，光本身是何种形态的客体，我们仍然是不得而知的，我们知道的一切只不过是光表现出的外在属性，即波粒二象性，这就足够了，至于光到底是何种存在，波还是粒子亦或是其他的某种什么客体，这都不是很重要。总之，这样我们就解决了“光怎么能同时是两种不同的事物”的诘难了，正确的认识“光的波粒二象性”，就要认识到波粒二象性是对属性而不是客体的描述。\goodbreak

光是波粒二象性的，描述光的物理量就可以相应的分为两组
\begin{itemize}
    \item 反映光的波动性的物理量主要包含：光速$c$、光的波长$\lambda$、光的频率$\nu$。
    \item 反映光的粒子性的物理量主要包含：质量$m$、动量$p$，能量$E$。
\end{itemize}
这两者间的桥梁就是\xref{fml:光子的能量}给出的$E=h\nu$，但是我们可以建立更丰富的联系。

\begin{BoxFormula}[普朗克--爱因斯坦关系式]
    光子的质量可以用频率或波长表示为
    \begin{Equation}&[质量]
        m=\frac{h\nu}{c^2}=\frac{h}{c\lambda}
    \end{Equation}
    光子的动量可以用频率或波长表示为
    \begin{Equation}&[动量]
        p=\frac{h\nu}{c}=\frac{h}{\lambda}
    \end{Equation}
    以上两式是描述光的波粒二象性的两个基本关系，通常合称为\uwave{普朗克--爱因斯坦关系式}。
\end{BoxFormula}

\begin{Proof}
    根据\fancyref{fml:爱因斯坦质能公式}
    \begin{Equation}
        E=mc^2
    \end{Equation}
    稍作变换，代入\fancyref{fml:光子的能量}，得到
    \begin{Equation}&[1]
        m=\frac{E}{c^2}=\frac{h\nu}{c}
    \end{Equation}
    我们知道波长和频率的关系是$\lambda=cT=c/\nu$，因此有
    \begin{Equation}&[2]
        m=\frac{h}{c\lambda}
    \end{Equation}
    而动量的公式就很容易理解了，我们知道
    \begin{Equation}
        p=mc
    \end{Equation}
    在上式中代入\xrefpeq{1}和\xrefpeq{2}即得\xrefpeq{动量}。
\end{Proof}
